\documentclass[conference]{IEEEtran}

\usepackage{iftex}
\ifPDFTeX
\usepackage[utf8]{inputenc}
\usepackage[T1]{fontenc}
\fi
\usepackage{amsthm}
\usepackage{amsmath,bm,amssymb}
\usepackage{newtxtext,newtxmath}
\usepackage{textcomp}
\interdisplaylinepenalty=2500
\usepackage{xcolor}
\usepackage{graphicx}
\usepackage{array,booktabs,makecell,adjustbox}
\usepackage{mathtools}
\usepackage{enumitem}
\usepackage{siunitx}
\usepackage[final]{microtype}
\usepackage{cite}
\usepackage[breaklinks=true,hidelinks]{hyperref}
\usepackage{bookmark}
\usepackage[nameinlink,capitalise,noabbrev]{cleveref}

% Figures are drawn into plots/kalman_carma by draw_plots.sh from the
% committed study data and moved here by CI.  Wrap every include so the
% manuscript compiles from a clean checkout.
\newcommand{\IncludeFig}[2][\linewidth]{%
  \IfFileExists{#2}{\includegraphics[width=#1]{#2}}%
                   {\fbox{\parbox{0.92\linewidth}{\centering\small
                     Figure \texttt{\detokenize{#2}} not generated.\\
                     Run \texttt{plots/kalman\_carma/draw\_plots.sh}.}}}%
}
\newcommand{\InputPart}[1]{%
  \IfFileExists{#1}{\input{#1}}{\fbox{\small Table \texttt{\detokenize{#1}} not generated.}}}

\DeclareRobustCommand{\vct}[1]{\bm{#1}}
\DeclareRobustCommand{\mat}[1]{\bm{#1}}
\newcommand{\R}{\mathbb{R}}
\newcommand{\C}{\mathbb{C}}
\newcommand{\Hr}{^{\mathsf H}}
\sisetup{detect-weight=true,detect-family=true,per-mode=symbol}
\emergencystretch=2em
\allowdisplaybreaks

\newtheoremstyle{bodysmall}{3pt}{3pt}{\normalfont\small}{0pt}{\bfseries\small}{.}{0.5em}{}
\theoremstyle{bodysmall}
\newtheorem{proposition}{Proposition}
\newtheorem{remark}{Remark}

\title{Coupled Three-Axis CARMA Priors for a Quaternion MEKF Marine
Motion Estimator: Correlated Excitation, Complexified Shared Modes and a
Constrained MCARMA (experimental draft)}
\author{%
  \IEEEauthorblockN{Mikhail S. Grushinskiy}
  \IEEEauthorblockA{Independent Researcher, USA, Fair Lawn\\
  Bareboat Necessities GitHub Project\\
  Email: mgrouch@users.github.com}%
}

\begin{document}
\maketitle

\begin{abstract}
The experimental CARMA3 quaternion multiplicative extended Kalman filter of
ocean-IMU models the wave displacement of each axis with an independent
CARMA(4,1) process. Linear waves couple the axes: surge and heave are in
quadrature, with a coherence near one in a single wave system. We compare
three coupled priors in the same filter: correlated excitation (21 states),
shared directional modes, a complexified CARMA(4,1) mode per axis driven by
Hermitian noise (33 states), and a constrained multivariate CARMA with a
heave response amplitude operator (35 states), all with a coupling estimated
from the IMU alone. In a paired Monte Carlo study (50 seeds, ten wave seas,
complete INS, simulation only) the shared modes lower the pooled
displacement error by 11.4\,\%, 4.8\,\% and 12.2\,\% (north, east, heave;
95\,\% intervals within $\pm0.6$ points), within 1.4 points of the oracle
coupling and with calm water unchanged. Correlated excitation gains
nothing, since a real coupling carries at most 41\,\% of the surge--heave
cross-spectral energy. The constrained MCARMA gains more (26.5\,\% heave),
but only with the benchmark hull's response operator; a 27\,\% frequency
error turns the gain into a loss of up to 25\,\%, and its worst update costs
\SI{824}{\micro s} against \SI{35}{\micro s}. A reversed coupling costs up to
38\,\%, so a deployable filter needs a coupling confidence and fallback. We
recommend the shared directional modes as a separate experimental filter;
the production CARMA3 filter is unchanged.
\end{abstract}

\section{Introduction and Status}
The experimental CARMA3 estimator of ocean-IMU is a quaternion
multiplicative extended Kalman filter (QMEKF) whose translational prior is one
CARMA(4,1) displacement process per axis (surge, sway, heave). The axes share
the measured operating point but their stochastic dynamics are independent.
Linear wave theory says they are not independent: one wave component moves a
floating hull along its direction of travel a quarter period behind the
elevation and heaves it through the heave response amplitude operator (RAO).
This paper asks whether a physically informed, coupled three-dimensional
CARMA prior can exploit these relations to improve QMEKF estimates of surge,
sway and heave beyond three independent CARMA(4,1) processes, and answers it
with exact analysis and paired Monte Carlo experiments in which coupling is
beneficial, irrelevant or harmful.

Three model families of increasing generality are compared with the
unchanged baseline: (A) correlated white-noise excitation through the
baseline's identical per-axis dynamics; (B) shared directional modes built
from a \emph{complexified} CARMA(4,1) whose real part is exactly the
baseline process and whose imaginary part is its quadrature companion;
(C) a constrained multivariate CARMA (MCARMA) with the heave RAO in its
transfer matrix. A and the modal form of the baseline are evaluated here as
filters; B is the standalone shared-modes filter, in its six-state-per-axis
minimal realization, with its own articles; C is analysed only. All are
experimental, run behind the unchanged wrapper,
startup, magnetic acquisition and tuner, and leave every default, the
shipping OU-III filter and the two-frame (TFG) variants untouched. Nothing
here is hardware evidence or a stability statement.

\section{The Baseline as Implemented}
\label{sec:baseline}
The error state is
$[\delta\vct\theta,\vct b_g,\vct x_x,\vct x_y,\vct x_z,\vct b_a]\in\R^{21}$
(NED, $z$ down) with one CARMA(4,1) block
$\vct x_i=[p_i,v_i,a_i,l_i]^{\mathsf T}$ per axis, \emph{axis-major}, so the
motion block has $\mat A_L=\mat I_3\otimes\mat A_C$. The per-axis transfer
function is $H(s)=gB(s)/P(s)$,
$P(s)=(s^2+2\zeta_1\omega_1s+\omega_1^2)(s^2+2\zeta_2\omega_2s+\omega_2^2)$,
$B(s)=s+\omega_{z0}$; the JONSWAP fit C41J has
$f_1=\SI{0.306945}{Hz}$, $\zeta_1=0.104$, $f_2=\SI{0.815}{Hz}$,
$\zeta_2=0.478$, $\omega_{z0}=0$ (a zero at DC). The noise enters $\dot a$
only (relative degree three), so $p,v,a$ are states and every accelerometer
row is a unit selector rotated by $\mat R_{wb}$. $\mat F=e^{\mat Ah}$ and
$\mat Q=\int_0^h e^{\mat At}\mat W e^{\mat A^{\mathsf T}t}dt$ are computed
exactly (scaled Taylor series and doublings). The wrapper commits every
\SI{0.1}{s} a smoothed operating point: one frequency scale $\kappa$ for every
axis and isotropic acceleration amplitudes
$\sigma_x=\sigma_y=\sigma_z$; the covariance is reset to the stationary prior
only on entering Live. Horizontally only an optional real $x$--$y$
correlation existed; there was no $x$--$z$ or $y$--$z$ term.

\section{Cross-Axis Physics}
\label{sec:physics}
The benchmark truth is a random-phase directional sea through explicit RAOs.
For one component of azimuth $d$ the displacement phasor is
$\vct t=[-i\,\mathrm{su}_N,\,-i\,\mathrm{su}_E,\,H_z(\omega)]\eta$:
horizontal motion along the direction of travel in quadrature with the
elevation and hull-averaged, heave through a second-order RAO
($f_n=\SI{0.55}{Hz}$, $\zeta=0.35$). The cross-spectral matrix
$\mat S(f)=\sum_j S_\eta(f)w_j\vct t_j\vct t_j\Hr$ is Hermitian positive
semidefinite. At the JONSWAP peak the $x$--$z$ phase is \SI{-62}{\degree},
not \SI{-90}{\degree} (the RAO lags), it drifts with frequency as the RAO
phase does, and the energy-weighted coherence is 0.97. Only 35--47\,\% of the
$x$--$z$ and $y$--$z$ cross-spectral energy is real (co-spectrum) on
single-system seas, while a constant complex coefficient captures 89--94\,\%
(\cref{tab:repr}). Vessel response and water-particle motion differ, and no
universal RAO is assumed below.

\InputPart{carma-coupled-repr-generated.tex-part}

\section{Coupled Models}
\label{sec:models}
\subsection{Experiment A: correlated CARMA(4,1)}
With identical per-axis dynamics and a real positive semidefinite
acceleration covariance $\mat\Sigma_w$,
\begin{equation}
\mat F=\mat I_3\otimes\mat F_C,\quad
\mat Q=\mat\Sigma_w\otimes\mat Q_C,\quad
\mat P_\infty=\mat\Sigma_w\otimes\mat P_C,
\label{eq:kron}
\end{equation}
$\mat Q_C,\mat P_C$ those of the unit-acceleration-variance C41J; the factor
order follows the axis-major state. With $\mat\Sigma_w$ diagonal this is the
baseline.
\begin{proposition}[Limit of correlated excitation]
For \eqref{eq:kron} the displacement cross-spectral matrix is
$\mat S_p(\omega)=|H(i\omega)|^2\mat\Sigma_w$: real at every frequency, with
cross-phase $0$ or $\pi$ and frequency-independent coherence
$\Sigma_{w,ij}^2/(\Sigma_{w,ii}\Sigma_{w,jj})$.
\end{proposition}
Hence A can carry only the co-spectrum, 35--47\,\% of the $x$--$z$ energy
of the truth.

\subsection{Experiment B: complexified shared modes}
A damped stochastic oscillator
$d\vct q=\left[\begin{smallmatrix}-\lambda&-\omega\\\omega&-\lambda\end{smallmatrix}\right]\vct q\,dt+\mat G\,d\vct W$
is the real form of a complex Ornstein--Uhlenbeck process
$\dot\zeta=\mu\zeta+\xi$, $\mu=-\lambda+i\omega$, whose spectrum sits at
positive frequency, so $\operatorname{Re}\zeta$ and $\operatorname{Im}\zeta$
are nearly a Hilbert pair \cite{gabor46,schreier10}. A single rotation has
relative degree one (velocity would be white). The following construction
raises it to three and keeps the baseline law exactly.
\begin{proposition}[Complexification]
\label{prop:cx}
Let $P$ be a real monic quartic with roots $p_1,\bar p_1,p_2,\bar p_2$,
$\operatorname{Im}p_k>0$, $P_+(s)=(s-p_1)(s-p_2)$, and let $D(s)=s^2+d_1s+d_0$
with $d_0=\sqrt\gamma$, $d_1=\sqrt{\beta+2d_0}$, where
$\tfrac12(|P_+(i\omega)|^2+|P_+(-i\omega)|^2)=\omega^4+\beta\omega^2+\gamma$.
Then $D$ is Hurwitz and, for every real-coefficient $B$,
\begin{equation}
\left|\tfrac{B}{P_+D}(i\omega)\right|^2+\left|\tfrac{B}{P_+D}(-i\omega)\right|^2
=2\left|\tfrac{B}{P}(i\omega)\right|^2 .
\label{eq:cx}
\end{equation}
Consequently, if $\zeta=(B/P_+D)(s)\,\xi$ with $\xi$ circular complex white
noise of intensity $2g^2$, the real part of $\zeta$ is a Gaussian process
with spectrum $g^2|B/P|^2$: the real CARMA process itself.
\end{proposition}
\begin{proof}
$|P_+(-i\omega)|=|P_-(i\omega)|$ with $P_-=\bar P_+$, so
$|P(i\omega)|^2=|P_+(i\omega)|^2|P_+(-i\omega)|^2$. The left side of
\eqref{eq:cx} is
$|B|^2(|P_+(-i\omega)|^2+|P_+(i\omega)|^2)/(|P|^2|D|^2)=2|B|^2/|P|^2$ by the
choice of $D$; the quartic in $\omega$ is even with leading coefficient one,
and $|D(i\omega)|^2=\omega^4+(d_1^2-2d_0)\omega^2+d_0^2$ matches it; $d_0,d_1>0$
because the left side is positive. For circular $\xi$ the real part has
spectrum $\tfrac14(S_\zeta(\omega)+S_\zeta(-\omega))$.
\end{proof}
The construction is exact for C41J (relative error \num{1.8e-14} of the
real-part spectrum, \num{7.7e-15} of its $[p,v,a]$ stationary covariance;
poles strictly stable over $\kappa\in[0.15,3.2]$). The single-sidedness is
not exact: the mirror ratio $|H(-i\omega)|^2/|H(i\omega)|^2$ is \num{8e-4} at
the peak, 0.6--2\,\% at 0.45--\SI{1}{Hz} and 6\,\% (17\,\%) at 0.15
(0.1)\,Hz for $\kappa=1$. In kinematic coordinates
$[\zeta,\dot\zeta,\ddot\zeta,l]$ the real state is
$[\operatorname{Re}(p,v,a,l),\operatorname{Im}(p,v,a,l)]$, the noise enters
$a$ only, and every readout remains a unit selector. A representation of
MCARMA processes as sums of complex-valued OU processes is given in
\cite{fasen21} (\cref{sec:modal}); Proposition~\ref{prop:cx} instead builds
the single-sided companion of a given real process.

Each axis carries such a mode; the axes share its dynamics and are driven by
circular noise with Hermitian positive semidefinite covariance $\mat\Sigma_c$:
\begin{equation}
\mat F=\mat I_3\otimes\mat F_8,\quad
\mat Q_{ab}=\operatorname{Re}(\Sigma_{c,ab})\mat Q_8
          +\operatorname{Im}(\Sigma_{c,ab})\mat J\mat Q_8,
\label{eq:hkron}
\end{equation}
$\mat J=\left[\begin{smallmatrix}\mat 0&-\mat I\\\mat I&\mat 0\end{smallmatrix}\right]$,
which commutes with $\mat F_8,\mat Q_8$. The displacement cross-spectrum is
$S_{ij}(\omega)=\tfrac{g^2}2[\Sigma_{c,ij}|H(i\omega)|^2+\bar\Sigma_{c,ij}|H(-i\omega)|^2]$:
$\operatorname{Re}\mat\Sigma_c$ is the co-spectrum, $\operatorname{Im}\mat\Sigma_c$
the quadrature spectrum. A real $\mat\Sigma_c$ reproduces Experiment A in law
(\num{6e-15}), a real diagonal one the baseline. Two of the eight real
states per axis are unobservable for every excitation; the filter carries the
law in its six-state minimal realization (27 error states), the standalone
shared-modes QMEKF of \texttt{src/kalman\_carma\_shared\_modes}. A
state-reduced variant (the mode on $z$ only,
$x,y=\operatorname{Re}(\beta\zeta_z)+$ residual C41J, 16 motion states)
attains the same steady-state errors within 0.1--1\,mm, but its coupling sits
in the output map and would move the estimate at every adaptation, so it was
not implemented.

\subsection{Complex-Valued Ornstein--Uhlenbeck Factorization of Multivariate CARMA Processes}
\label{sec:modal}
\emph{Theory.} Fasen-Hartmann and Scholz \cite{fasen21} decompose a
$d$-variate MCARMA$(p,q)$ process with autoregressive polynomial
$\mat P(z)=\mat Iz^p+\mat A_1z^{p-1}+\dots+\mat A_p$ into a sum of $p$
multivariate complex-valued OU processes,
$\vct Y=\sum_{k=1}^p\vct Y^{(k)}$,
$d\vct Y^{(k)}_t=\mat R_k\vct Y^{(k)}_tdt+\mat C_k\,d\vct L_t$, where
$\mat R_1,\dots,\mat R_p$ are right solvents of $\mat P$ and the input
matrices $\mat C_k$ are residues of the transfer function at them. It needs
an irreducible fraction and a complete set of regular right solvents (for
instance, distinct eigenvalues of the block companion matrix); with repeated
solvents the process is not such a sum. All components are driven by the
\emph{same} L\'evy process, so they are dependent, with cross-covariances
given by the residues, and the sampled process is a VARMA$(p,p-1)$. The
state-space characterization of MCARMA processes in \cite{schlemm12} was
corrected in 2024 \cite{stelzer24}. The correction restates it through the
observer and controller canonical forms belonging to left and right
matrix-fraction descriptions with monic denominators, and this paper uses
it in that form. Every per-axis polynomial here is scalar, so the left and
right descriptions coincide and the modal construction below is unaffected.
The transfer matrices of the three-axis models
($\mat H=\mat P^{-1}\mat B$ with diagonal $\mat P$, and
$\operatorname{diag}(1,1,R_z)\mat H_c$ for C) are written as left fractions.

\emph{The modal form of CARMA(4,1).} For $d=1$, $p=4$ the solvents are the
poles. C41J has two complex-conjugate pairs, and its scalar partial-fraction
form (\cite{brockwell11}, Prop.~2) is
\begin{equation}
p_t=2\operatorname{Re}(z_{1,t}+z_{2,t}),\quad
dz_k=\mu_kz_k\,dt+r_kg\,dW_t,\quad r_k=\frac{B(\mu_k)}{P'(\mu_k)},
\label{eq:modal}
\end{equation}
with $\mu_k$ the upper poles and one common real Brownian motion $W$. At
$\kappa=1$, $\mu_{1,2}=-0.2009+1.9181i,\ -2.4497+4.4991i$ (rad/s) and
$r_{1,2}=0.02078-0.00587i,\ -0.02078+0.01289i$. With the real state
$[\operatorname{Re}z_1,\operatorname{Im}z_1,\operatorname{Re}z_2,\operatorname{Im}z_2]$
per axis (12 motion states, as in the baseline), $\mat A$ is block-diagonal
with damped rotations
$\left[\begin{smallmatrix}\operatorname{Re}\mu_k&-\operatorname{Im}\mu_k\\
\operatorname{Im}\mu_k&\operatorname{Re}\mu_k\end{smallmatrix}\right]$, the
noise input is $\vct G=g[\operatorname{Re}r_1,\operatorname{Im}r_1,
\operatorname{Re}r_2,\operatorname{Im}r_2]^{\mathsf T}$, and $p,v,a$ are the
dense readouts $\vct c_p=[2,0,2,0]$, $\vct c_p\mat A$, $\vct c_p\mat A^2$.
$\mat F$ is closed form (rotation--decay blocks). So is $\mat Q$, from
$E[z_k\bar z_l]=g^2r_k\bar r_l(e^{(\mu_k+\bar\mu_l)h}-1)/(\mu_k+\bar\mu_l)$
and $E[z_kz_l]=g^2r_kr_l(e^{(\mu_k+\mu_l)h}-1)/(\mu_k+\mu_l)$, including
the cross-mode blocks that the common noise creates. The limit
$h\to\infty$ gives $\mat P_\infty$.
\begin{proposition}[Modal kinematics]
\label{prop:modal}
(i) $\sum_k2\operatorname{Re}(r_k\mu_k^j)$ is the $j$-th Markov parameter of
$B/P$, so $\sum\operatorname{Re}r_k=\sum\operatorname{Re}(r_k\mu_k)=0$ and
$\vct c_p\vct G=\vct c_p\mat A\vct G=0\neq\vct c_p\mat A^2\vct G$: relative
degree three, so $v=\dot p$ and $a=\ddot p$ are noise-free readouts.
(ii) Driving each mode by its own noise of the same intensity removes the
cross-mode covariance and the cancellations. (iii) A readout rotated by a
constant phase, $2\operatorname{Re}(e^{i\varphi}(z_1+z_2))$, has
$\vct c\vct G=-2g\sin\varphi\sum_k\operatorname{Im}r_k\neq0$
($\sum\operatorname{Im}r_k=0.00702$), so its relative degree is one.
(iv) Per-mode complex gains $\beta_k$ on an axis driven by the same noise
keep relative degree three iff
$\sum\operatorname{Re}(\beta_kr_k)=\sum\operatorname{Re}(\beta_kr_k\mu_k)=0$,
that is, iff the axis is $(b_1B(s)+b_0)/P(s)\,\xi$ with real $b_0,b_1$. The
only admissible cross-axis phase is then
$-\arctan\bigl(b_0/(b_1\omega)\bigr)$ (for $\omega_{z0}=0$). It is never
constant, and $b_0=0$ leaves real coupling only (Experiment A).
\end{proposition}
\begin{proof}
(i) follows from expanding $B/P=\sum_k[r_k/(s-\mu_k)+\bar r_k/(s-\bar\mu_k)]$ in
powers of $1/s$, since $\deg P-\deg B=3$. (iv) With distinct poles, the
gains $\beta_kr_k$ correspond one-to-one to real numerators $N$ with
$\deg N\le3$ over $P$, and the relative degree is three iff $\deg N\le1$,
i.e.\ $N=b_1B+b_0$; (iii) is the special case $\beta_k=e^{i\varphi}$.
\end{proof}
Numerically (\cref{tab:modal}): over $\kappa\in[0.2,2.5]$ the modal form
reproduces the C41J transfer function to \num{3e-14}, its stationary
covariance to \num{9e-15}, $\mat F$ and $\mat Q$ under the similarity
$\vct x_{\mathrm{base}}=\mat T\vct x_{\mathrm{modal}}$ to \num{2e-15} and
\num{7e-15}, and the displacement autocovariance at 0.5--\SI{5}{s} lags to
\num{7e-15}; $\operatorname{cond}\mat T\le196$. Independently driven
oscillators instead give $Q_{pp}\propto h$ (log-slope 1.00 against 4.98 for
the exact model): white velocity, 1.387 times the displacement variance, and
an acceleration spectrum that rises at high frequency (log$_2$ slope $+1.97$
per octave against $-2.04$), i.e.\ no finite acceleration variance. That is
the naive ``independent damped oscillator modes'' model, and it is rejected.
The complexification of \cref{prop:cx} is consistent with this theory. It
is not a decomposition of C41J but a complex CARMA with the upper poles
$\mu_k$ and the two real poles of $D$ ($-7.93$, $-1.25$\,rad/s), and its
relative degree is three by construction. That is why arbitrary phases need
six states per axis (its minimal realization) while the four-state modal form
admits only the family (iv).

\emph{Validity and numerics.} The form requires distinct (regular) poles.
The motion block refuses shapes whose two poles nearly coincide
($|\mu_1-\mu_2|<10^{-3}\max|\mu_k|$) and keeps the baseline realization as the
fallback. The modal coordinates depend on $\kappa$ ($v=\vct c_p\mat
A(\kappa)\vct x$), so at every model commit the filter applies the
re-coordinatization $\mat M=\mat T_{\mathrm{new}}^{-1}\mat T_{\mathrm{old}}$
to the state and, as $\mat S\mat P\mat S^{\mathsf T}$, to the covariance.
The baseline coordinates and the posterior are therefore continuous. In
double precision the modal QMEKF equals the baseline to \num{2.7e-13}
(fixed model) and \num{7.5e-13} (\SI{10}{Hz} changes of $\kappa$ and
amplitude); without the re-coordinatization it does not. Its simulated
$p,v,a$ variances match $\mat P_\infty$ to 2.1\,\% (11\,h at \SI{50}{Hz}).

\emph{QMEKF implementation.} The state ordering, the attitude, bias,
magnetometer, reset and Joseph code are the baseline's. The modal blocks
replace the CARMA blocks (axis-major, $\mat F=\mat I_3\otimes\mat F_4$,
$\mat Q=\mat\Sigma\otimes\mat Q_4$), and the accelerometer Jacobian columns
become $\mat R_{wb}^{\mathsf T}\vct e_i\,\vct c_a$ per axis (four terms
instead of one). The attitude/bias--motion cross-covariances propagate
through the same congruence, and the multiplicative attitude reset is
unchanged because it does not touch motion states. Position and velocity,
with their covariance, are read through the dense rows. Runtime arithmetic
is real; complex arithmetic in double is used only in the \SI{10}{Hz} model
rebuild (poles, residues, $\mat T$, closed-form $\mat F,\mat Q$).

\emph{Directional modal coupling.} By \cref{prop:modal}(iv), the only
kinematically admissible directional coupling of the four-state modal form
gives each horizontal axis a numerator zero, $(B+b_{0,i})/P$, with a real
excitation between axes. The deployable variant (Form \emph{Numerator})
takes $b_{0,i}=-\omega_{op}\tan(\arg S_{iz})\rho_{iz}$ ($|\arg|\le
\SI{70}{\degree}$, $\rho$ the excitation coherence) from the same Hermitian
band estimate as B. Its excitation is
$\operatorname{Re}(\mat D\mat\Sigma_c\mat D\Hr)$,
$\mat D=\operatorname{diag}(1/u_x,1/u_y,1)$, $u_i=1-ib_{0,i}/\omega_{op}$.
This matches the band gain and the phase $\arctan(\rho\tan\arg S_{iz})$ at
$f_{op}$ to \num{4e-16} (test [17]) and keeps the excitation positive
semidefinite. The price is a nonzero displacement density at DC on the
coupled axes. In the steady state (\cref{tab:steady}, ``Modal dir.'')
heave improves by about as much as with B (JONSWAP 10.2$\to$9.3\,mm), but
the horizontal error grows (7.4$\to$9.7\,mm) from that low-frequency
density; on the mixed sea the trade reverses horizontally
(89.7$\to$68.6\,mm). With an exact representation it is the experimental
candidate of this section. The modal form with a real excitation
(\texttt{MOD:est}) is Experiment A in law and gives identical Monte Carlo
rows.

\InputPart{carma-coupled-modal-generated.tex-part}

\subsection{Experiment C: constrained MCARMA}
The general MCARMA(4,1) $\mat P(s)=s^4\mat I+\mat A_1s^3+\dots+\mat A_4$,
$\mat B(s)=\mat B_0s+\mat B_1$, $\mat S_p=\mat H\mat\Sigma\mat H^*$,
$\mat H=\mat P^{-1}\mat B$ \cite{marquardt07,schlemm12,stelzer24} has 54
coefficients. Its Whittle Fisher information \cite{whittle53,fasen22} from
\SI{20}{min} of three-axis acceleration at a generic point near the baseline
has rank 51 (the three null directions are the $\mat B\mapsto\mat B\mat U$
gauge), but 44\,\% of the coefficients have a relative Cram\'er--Rao bound
above 100\,\% (off-diagonal autoregressive coefficients 60--115\,\%,
$\mat B_1$ essentially unidentified). The constrained form keeps
the physics and 12 identifiable parameters: complexified modes with Hermitian
excitation and the vertical output through the heave RAO,
\begin{equation}
\mat H(s)=\operatorname{diag}\bigl(1,1,R_z(s)\bigr)\mat H_c(s),\quad
R_z(s)=\frac{\omega_n^2}{s^2+2\zeta\omega_ns+\omega_n^2},
\end{equation}
so that the $x$--$z$ phase follows $\arg\Sigma_{xz}-\arg R_z(i\omega)$
across frequency (\SI{34}{\degree} between 0.75 and $1.3f_{op}$ for the
benchmark hull). The RAO parameters are identifiable (relative CRLB
0.8\,\% and 2.9\,\% in \SI{20}{min}). The axes no longer share dynamics, so a
filter would need the exact discretization of a joint 26-state system and
hull knowledge (the RAO). The form is analysed (\cref{sec:steady}) but not
implemented.

\subsection{Two modes, swell and wind sea, crossing seas}
Two complexified modes per axis (48 motion states) with oracle frequencies
cut the steady-state error on the mixed sea from 89.7/93.6\,mm
(horizontal/heave) to 27.6/29.7\,mm, but 29.4/36.2\,mm of that comes from
the two-peak marginal shape alone (no coupling); on two crossing systems of
similar frequency coupling gains nothing. No second-peak frequency
estimator exists in the wrapper, so this variant was analysed, not
implemented.

\section{Coupling Estimation and Frames}
\label{sec:estimation}
The filter's own world-frame acceleration
$\vct a=\mat R_{wb}^{\mathsf T}(\vct f-\hat{\vct b}_a)+g\vct e_z$ is
complex-demodulated at the operating frequency \cite{bingham67}: a two-pole
envelope low-pass at $0.15f_{op}$ and a \SI{240}{s} exponential average of
$\vct A\vct A\Hr$ give a Hermitian positive semidefinite band cross-spectrum
$\hat{\mat S}$. Because $\vct a=-\omega^2\vct p$ on every axis, its ratios
are the displacement regressions at $f_{op}$; the quadrature sign carries
the direction of travel, removing the \SI{180}{\degree} ambiguity of
co-spectral estimates \cite{longuethiggins63,kuik88}. The excitation keeps
the baseline diagonal, the wrapper's isotropic $\sigma_h^2\mat I$ horizontally
(sensor noise, tilt and bias leakage in the horizontal band are not wave
energy) and $\sigma_z^2$, and adds only the cross terms
\begin{equation}
\vct b=s\,c\,\hat{\mat S}_{hz},\quad c=\frac{|R_z|^2\sigma_z^2}{\hat S_{zz}},\quad
s=\mathrm{gain}\,\max\Bigl(0,1-\frac1{N_{\mathrm{eff}}\gamma^2}\Bigr),
\label{eq:mapping}
\end{equation}
so every regression of a horizontal axis on $z$ at $f_{op}$ (co- and
quadrature spectrum) is the measured one. Here $\gamma^2$ is the debiased
multiple coherence of $z$ on $(x,y)$ and $N_{\mathrm{eff}}$ the number of
independent envelope samples. $\vct b$ is scaled down if
$\vct b\Hr\mat H^{-1}\vct b$ would exceed $0.999\sigma_z^2$, and
$\mat\Sigma=\mat R^{-1}\mat M\mat R^{-\mathsf H}$ for a model whose axes
respond with $R_i$ at $f_{op}$ (the identity for every implemented model). $\mat\Sigma$ is positive
semidefinite by construction and reduces to the baseline diagonal without
coupling evidence. It is also equivariant under rotations about the vertical
(to \num{7e-16}), since every factor is an invariant, so the prior does not
depend on how the waves sit relative to north. A form that also scales the
\emph{measured} horizontal block by $c$ reads horizontal tilt and bias
leakage as wave energy and inflates the calm-water error; it is an ablation
option. The band fraction (0.15) and horizon (\SI{240}{s}) are the
estimator's defaults. ORACLE runs feed the same
estimator with the true wave acceleration; deployable runs use the IMU only.
The world frame is the physically right one: a wave field is fixed in the
world, so the coupling survives heading changes, while a coupling held in
the keel frame turns with the hull.

\section{Implementation and Tests}
\label{sec:impl}
The baseline class gained one defaulted template parameter
\texttt{Motion = void}; the default instantiation compiles to the unchanged
filter. This was checked bit for bit on every column of paired benchmark
rows (six CARMA variants, jitter, turns, level~1; re-run after the last core
change) and on the ESP32-S3 object code of the complete wrapper, which is
instruction- and literal-identical to the original apart from symbol names.
A coupled filter shares the attitude, bias, magnetometer, reset and Joseph
code line for line and is used as the wrapper's inner filter, so the wrapper
and every default are untouched. Only the linear block, its readouts (unit
selectors or dense rows), the accelerometer columns and the modal
re-coordinatization dispatch. A 54-check numerical test suite covers
reduction to the baseline (double precision: \num{1e-14}), positive
semidefiniteness, discretization accuracy, cross-covariance propagation,
cross-spectral phase, frame consistency, attitude resets, changing and
near-zero excitation, sensor-noise-only input, long runs and the modal form
(\cref{sec:modal}). The shared directional modes, model-matched consistency
and wrong-coupling degradation included, are tested with their filter.

\section{Steady-State Analysis}
\label{sec:steady}
\Cref{tab:steady} gives the exact steady-state (200\,Hz DARE) error of every
model against the \emph{true} cross-spectral matrix, white accelerometer
noise only, oracle parameters. It is a non-promoting feasibility diagnostic:
it has no tilt, bias or adaptation.

\InputPart{carma-coupled-steady-generated.tex-part}

\section{Monte Carlo Protocol}
\label{sec:protocol}
Every row runs one scenario and seed: one directional sea through the hull
RAOs, one IMU/magnetometer stream (accelerometer \SI{1.5}{mg} white, turn-on
bias up to \SI{8}{mg} with random walk; gyroscope \SI{0.0016}{rad/s} white,
bias up to \SI{0.05}{\degree/s}; magnetometer \SI{0.8}{\micro T} at
\SI{25}{Hz} with hard/soft iron, and in one scenario periodic
\SI{15}{\micro T} disturbances with scale and misalignment residuals), fed
identically to every filter (paired). Errors are taken against the wave
coordinate from \SI{300}{s} to \SI{1200}{s}. The estimator's hyper-parameters
are its defaults; the final comparison uses 50 seeds, the ablations 30. Changes are
$1-\overline{\mathrm{RMS}}_{\mathrm{cand}}/\overline{\mathrm{RMS}}_{\mathrm{base}}$
with 95\,\% bootstrap intervals over seeds.

\section{Results}
\label{sec:results}
\InputPart{carma-coupled-results-generated.tex-part}

\begin{figure*}[t]
\centering
\IncludeFig[0.98\textwidth]{rms_change_by_sea.png}
\caption{Paired change of the RMS displacement error against the baseline,
per axis and sea, 50 seeds (positive = better), 95\,\% bootstrap
intervals.}
\label{fig:rms}
\end{figure*}

\begin{figure*}[t]
\centering
\IncludeFig[0.98\textwidth]{spectra_mixed_swell_wind.png}
\caption{Mixed swell and wind sea: displacement PSDs, error PSDs,
cross-spectral magnitudes, phases and coherences of the truth and of every
estimate (Welch, 10 seeds).}
\label{fig:spec}
\end{figure*}

\begin{figure}[t]
\centering
\IncludeFig{cross_phase_models.png}
\caption{Cross-spectral phase of the oracle-parameterized models against the
truth: A is real, B holds a constant phase, C follows the RAO.}
\label{fig:phase}
\end{figure}

\begin{figure}[t]
\centering
\IncludeFig{consistency_nees_nis.png}
\caption{NEES of $[\vct p,\vct v]$ and accelerometer NIS per sea and model.}
\label{fig:nees}
\end{figure}

\begin{figure}[t]
\centering
\IncludeFig{ablation_sensitivity.png}
\caption{Wrong, frozen and keel-frame coupling and RAO mismatch (30 seeds).}
\label{fig:abl}
\end{figure}

\begin{figure}[t]
\centering
\IncludeFig{direction_change.png}
\caption{Wave direction turning \SI{90}{\degree}: rolling error and the
applied $x$--$z$ phase.}
\label{fig:dir}
\end{figure}

\begin{figure}[t]
\centering
\IncludeFig{heading_change.png}
\caption{\SI{90}{\degree} zigzag heading changes.}
\label{fig:turns}
\end{figure}

\begin{figure}[t]
\centering
\IncludeFig{onset_cessation.png}
\caption{Wave onset at \SI{400}{s} and cessation at \SI{900}{s}.}
\label{fig:onoff}
\end{figure}

\begin{figure}[t]
\centering
\IncludeFig{cost_time_memory.png}
\caption{Host update time and ESP32-S3 code, object and stack size.}
\label{fig:cost}
\end{figure}

\section{Discussion and Recommendation}
\label{sec:discussion}
\textbf{What the gain is.} The shared-mode prior adds one statistic to the
independent-axis model: the Hermitian band cross-spectrum, whose imaginary
part is the quadrature relation of linear waves. A real (co-spectrum)
coupling represents 28--41\,\% of the truth's surge--heave cross-spectral
energy, a constant phase 84--94\,\%; correspondingly the correlated prior
gains nothing (heave $+0.6$\,\%, horizontal $-1.3$ and $-3.3$\,\%) while
the shared modes gain on every wave sea except short-period wind waves
(heave $-1.2$\,\%), most on long-period swell ($+23.7$\,\% heave) and on
two-system seas ($+16.1$\,\%). The estimated coupling is within 1.4 points
of the oracle, so the IMU-only estimator is close to the information the
model can use.

\textbf{What it costs.} Twice the host update time (\SI{7.0}{\micro s}
against \SI{3.6}{\micro s} mean), a \SI{15}{kB} filter object on the
ESP32-S3, and fragility: a reversed coupling costs 8--38\,\% and a
\SI{90}{\degree} error 6--25\,\%. The coupling belongs in the world frame (a
keel-frame coupling loses the gain through heading changes), and its
estimate must be shrunk and monitored with a fallback to independent axes.

\textbf{Why not the MCARMA.} The unrestricted MCARMA(4,1) is not identifiable
from one IMU (Fisher rank 51 of 54). The constrained form with a heave
response operator is the most accurate model here, but part of its gain is
hull knowledge: its operator prior alone gains 2--7\,\%, it uses the
benchmark hull's operator, and a 27\,\% natural-frequency error costs up to
25\,\%. It also doubles the accelerometer NIS, and its worst update
(\SI{824}{\micro s}, 16\,\% of the \SI{5}{ms} step on the host) is
unsuited to the AtomS3R.

\textbf{Recommendation (Option C).} Adopt the shared directional modes as a
separate experimental filter with a coupling confidence and fallback, not
as a change to the production CARMA3 filter, which is unchanged. That filter
is \texttt{KalmanCarmaSharedModes} (\texttt{docs/kalman-carma-shared-modes.md}).
All results are simulation; no hardware evidence is claimed.

\section*{Reproduction}
\texttt{make -C tests/kalman\_carma build test}; then
\texttt{python3 tools/carma3\_coupled\_analysis.py} and
\texttt{python3 tools/carma3\_coupled\_benchmark.py}
\texttt{run}, \texttt{ablate}, \texttt{spectra},
\texttt{traces}, \texttt{timing}, \texttt{report}. Raw rows, spectra, traces
and generated tables are in \texttt{reports/results/carma3\_coupled}; the
ESP32-S3 sizes come from
\texttt{tests/kalman\_carma/arduino/carma3\_size\_probe/build.sh} with the
Espressif \texttt{esp-14.2.0\_20241119} toolchain.

\begin{thebibliography}{99}
\bibitem{brockwell01}
P. J. Brockwell, ``L\'evy-driven CARMA processes,'' \emph{Ann. Inst. Statist.
Math.}, vol. 53, pp. 113--124, 2001.
\bibitem{marquardt07}
T. Marquardt and R. Stelzer, ``Multivariate CARMA processes,''
\emph{Stochastic Process. Appl.}, vol. 117, no. 1, pp. 96--120, 2007.
\bibitem{schlemm12}
E. Schlemm and R. Stelzer, ``Multivariate CARMA processes, continuous-time
state space models and complete regularity of the innovations of the sampled
processes,'' \emph{Bernoulli}, vol. 18, no. 1, pp. 46--63, 2012.
\bibitem{stelzer24}
R. Stelzer, ``Correction to: Multivariate CARMA processes, continuous-time
state space models and complete regularity of the innovations of the sampled
processes,'' arXiv:2411.06935, 2024.
\bibitem{schlemm12qml}
E. Schlemm and R. Stelzer, ``Quasi maximum likelihood estimation for strongly
mixing state space models and multivariate L\'evy-driven CARMA processes,''
\emph{Electron. J. Stat.}, vol. 6, pp. 2185--2234, 2012.
\bibitem{fasen17}
V. Fasen and S. Kimmig, ``Information criteria for multivariate CARMA
processes,'' \emph{Bernoulli}, vol. 23, no. 4A, pp. 2860--2886, 2017.
\bibitem{fasen22}
V. Fasen-Hartmann and C. Mayer, ``Whittle estimation for continuous-time
stationary state space models with finite second moments,'' \emph{Ann. Inst.
Statist. Math.}, vol. 74, pp. 233--270, 2022.
\bibitem{fasen21}
V. Fasen-Hartmann and M. Scholz, ``Factorization and discrete-time
representation of multivariate CARMA processes,'' arXiv:2102.11681, 2021.
\bibitem{brockwell11}
P. J. Brockwell, R. A. Davis and Y. Yang, ``Estimation for non-negative
L\'evy-driven CARMA processes,'' \emph{J. Bus. Econom. Statist.}, vol. 29,
pp. 250--259, 2011.
\bibitem{whittle53}
P. Whittle, ``The analysis of multiple stationary time series,'' \emph{J. R.
Stat. Soc. B}, vol. 15, no. 1, pp. 125--139, 1953.
\bibitem{vanloan78}
C. F. Van Loan, ``Computing integrals involving the matrix exponential,''
\emph{IEEE Trans. Autom. Control}, vol. 23, no. 3, pp. 395--404, 1978.
\bibitem{gabor46}
D. Gabor, ``Theory of communication,'' \emph{J. IEE}, vol. 93, pt. III,
no. 26, pp. 429--457, 1946.
\bibitem{neeser93}
F. D. Neeser and J. L. Massey, ``Proper complex random processes with
applications to information theory,'' \emph{IEEE Trans. Inf. Theory},
vol. 39, no. 4, pp. 1293--1302, 1993.
\bibitem{schreier10}
P. J. Schreier and L. L. Scharf, \emph{Statistical Signal Processing of
Complex-Valued Data}. Cambridge University Press, 2010.
\bibitem{bingham67}
C. Bingham, M. D. Godfrey and J. W. Tukey, ``Modern techniques of power
spectrum estimation,'' \emph{IEEE Trans. Audio Electroacoust.}, vol. 15,
no. 2, pp. 56--66, 1967.
\bibitem{hasselmann73}
K. Hasselmann \emph{et al.}, ``Measurements of wind-wave growth and swell
decay during the Joint North Sea Wave Project (JONSWAP),''
\emph{Erg\"anzungsh. Dtsch. Hydrogr. Z., Reihe A}, no. 12, 1973.
\bibitem{pierson64}
W. J. Pierson and L. Moskowitz, ``A proposed spectral form for fully
developed wind seas based on the similarity theory of S. A.
Kitaigorodskii,'' \emph{J. Geophys. Res.}, vol. 69, no. 24,
pp. 5181--5190, 1964.
\bibitem{mitsuyasu75}
H. Mitsuyasu \emph{et al.}, ``Observations of the directional spectrum of
ocean waves using a cloverleaf buoy,'' \emph{J. Phys. Oceanogr.}, vol. 5,
no. 4, pp. 750--760, 1975.
\bibitem{longuethiggins63}
M. S. Longuet-Higgins, D. E. Cartwright and N. D. Smith, ``Observations of the
directional spectrum of sea waves using the motions of a floating buoy,'' in
\emph{Ocean Wave Spectra}. Prentice-Hall, 1963, pp. 111--136.
\bibitem{kuik88}
A. J. Kuik, G. Ph. van Vledder and L. H. Holthuijsen, ``A method for the
routine analysis of pitch-and-roll buoy wave data,'' \emph{J. Phys.
Oceanogr.}, vol. 18, no. 7, pp. 1020--1034, 1988.
\bibitem{stdenis53}
M. St.~Denis and W. J. Pierson, ``On the motions of ships in confused
seas,'' \emph{Trans. SNAME}, vol. 61, pp. 280--357, 1953.
\bibitem{faltinsen90}
O. M. Faltinsen, \emph{Sea Loads on Ships and Offshore Structures}.
Cambridge University Press, 1990.
\bibitem{fossen11}
T. I. Fossen, \emph{Handbook of Marine Craft Hydrodynamics and Motion
Control}. Wiley, 2011.
\bibitem{nielsen06}
U. D. Nielsen, ``Estimations of on-site directional wave spectra from
measured ship responses,'' \emph{Marine Structures}, vol. 19, no. 1,
pp. 33--69, 2006.
\bibitem{tannuri03}
E. A. Tannuri, J. V. Sparano, A. N. Simos and J. J. Da Cruz, ``Estimating
directional wave spectrum based on stationary ship motion measurements,''
\emph{Appl. Ocean Res.}, vol. 25, no. 5, pp. 243--261, 2003.
\bibitem{saelid83}
S. S\ae lid, N. A. Jenssen and J. G. Balchen, ``Design and analysis of a
dynamic positioning system based on Kalman filtering and optimal control,''
\emph{IEEE Trans. Autom. Control}, vol. 28, no. 3, pp. 331--339, 1983.
\bibitem{fossen09}
T. I. Fossen and T. Perez, ``Kalman filtering for positioning and heading
control of ships and offshore rigs,'' \emph{IEEE Control Syst. Mag.},
vol. 29, no. 6, pp. 32--46, 2009.
\bibitem{markley03}
F. L. Markley, ``Attitude error representations for Kalman filtering,''
\emph{J. Guid. Control Dyn.}, vol. 26, no. 2, pp. 311--317, 2003.
\bibitem{barshalom01}
Y. Bar-Shalom, X. R. Li and T. Kirubarajan, \emph{Estimation with
Applications to Tracking and Navigation}. Wiley, 2001.
\end{thebibliography}

\end{document}
